Vegna þess að \(AB = AC\) er línan \(AD\) samhverfuás þríhyrningsins og þar með einnig miðlína og hæð á grunninn \(BC\).  
Því er \(D\) miðpunktur \(BC\) og \(BD = DC\).

Skoðum rétthyrnda þríhyrninginn \(ABD\): Hér er \(AB = 10\) (langhlið) og \(AD = 8\). Þá fæst
\[
BD = \sqrt{AB^{2} - AD^{2}} = \sqrt{10^{2} - 8^{2}} = \sqrt{100 - 64} = \sqrt{36} = 6.
\]
Þar sem \(D\) er miðpunktur \(BC\), þá er
\[
BC = 2 \cdot BD = 2 \cdot 6 = 12.
\]
